% Labels of the draft are printed as \dref{label} (texttt), since this note is compiled on its own.
\providecommand{\dref}[1]{\texttt{\detokenize{#1}}}
\section{Referee 3, Majors 3 and 4: $\beta_0$ identification, grouping, scales and scheme}\label{sec:ref3beta}
% Scripts: notes/scripts/ref3_beta/b1_vertices.py (+ b1.out), b2_coefficients.py (+ b2.out).
\noindent\textbf{Scope (author's instruction).} The calculation and the result of the paper are for pure Yang--Mills theory, $N_f=0$; nothing in this note changes that, and no $N_f$ term is added to the result. Quark loops are used only in a clearly separated remark (Sec.~\ref{sec:b-blm}), as a BLM-type device that answers the referee's question which of the numbers $\frac{11}{6}$ are $\beta_0/(2N_c)$ and which are integrals of the splitting function. Items (c)--(e) are for $N_f=0$.

\noindent\textbf{Results in brief.}
(1) \emph{Identification} (Sec.~\ref{sec:b-blm}): the $\frac{11}{3}$ of Row~I.II, the $-\frac{11}{6}$ of $N_{2b\,\rm UV}$ of Row~I.V, the $\frac{11}{6}$ of the finite kernel $\mathfrak K^{\rm V}_{2b}$ (Eqs.~\dref{eq:sGIkernels}, \dref{eq:app_V2}), the $\frac{11}{12}$ of the normalisation, of the counterterm and of the RC term, and the $\frac{\beta_0}{2}$ of $\delta(1-\xi)$ all become $\beta_0$-type (they are $\int d\xi$ of the \emph{complete} one-loop gluon splitting kernel, virtual or UV); the $-\frac{11}{6}$ of the large-distance tail, Eq.~(\dref{eq:dg_tailtot}), is the double-plus integral of the \emph{real} $\mathcal P_{gg+}$ and stays $\frac{11}{6}$ for any $N_f$.
(2) \emph{Referee's error item} (Sec.~\ref{sec:b-err}): no contradiction; the $\frac{11}{6}$ of $\mathfrak K^{\rm V}_{2b}$ is the integrand of the subtraction term whose dimensionally regularised $z$ integral cancels Row~I.II; the text has to say so.
(3) \emph{Grouping} (Sec.~\ref{sec:b-group}): the cancellation of the poles and of $\ln\mu^2$ between Rows~I.II and I.V is a result; ``the emission coupling does not run with $\v x-\v w$'' is a statement about a grouping. With the pair's colour evaluated at its centre of mass, the self-energy part of Row~I.II, \emph{including its constant $\frac{67}{9}-\frac{\pi^2}{3}$}, is cancelled identically (exact LCPT identity) by the unresolved part of Row~I.V, and the $\beta_0$-type coefficient survives only in the region where the pair resolves the target. This bears on the constant (Sec.~\ref{sec:b-const}).
(4) \emph{Scales} (Sec.~\ref{sec:b-scales}): with $\mu_R\neq\mu_F$ the result contains $+\frac{\alpha_s\beta_0}{4\pi}\ln\frac{\mu_R^2}{\mu_F^2}\times$LO; formulas given.
(5) \emph{Scheme} (Sec.~\ref{sec:b-scheme}): the paper is consistently CDR ($d_\perp=2-\epsilon$ polarisations wherever a pole multiplies a polarisation trace); with two-dimensional polarisations the Row~I.II constant becomes $\frac{64}{9}-\frac{\pi^2}{3}$ ($K\to K-\frac{N_c}{6}$), Row~I.V shifts by the opposite amount, Row~II.IV by $-\mathcal C\int d\xi\,2N_c\xi(1-\xi)\Pi(\xi)$, and $\alpha_s^{\rm DR}=\alpha_s^{\MS}(1+\frac{\alpha_sN_c}{12\pi})$; these close.

\subsection{Remark: BLM-type identification with a quark loop (not part of the result)}\label{sec:b-blm}

\paragraph{Vertices (\texttt{b1}).} The bracket of the paper's triple-gluon matrix element, Eq.~(\dref{eq:m3}), is \emph{exactly} minus the covariant vertex $\Gamma^{\mu\nu\rho}$ contracted with light-cone polarisation vectors ($\epsilon^+=0$, $\epsilon^-=\boldsymbol\epsilon\cdot\v k/k^+$), at random kinematics (ratio $-1$ to all digits, \texttt{b1}~[1]); the paper's vertex has the standard normalisation $gf^{abc}\Gamma/[(2\pi)^{3/2}\sqrt{2k^+2p^+2q^+}]$. With the quark vertex $g\,t^a\bar u\slashed\epsilon v$ in the same normalisation (explicit light-cone helicity spinors), the angular-averaged squared vertices satisfy, for every $\xi$ and also at non-zero parent momentum,
\begin{align}
\frac{n_fT_R\sum|V_{q\bar q}|^2}{\frac12N_c\sum|V_{gg}|^2}=\frac{n_fT_R}{N_c}\,\frac{\xi^2+(1-\xi)^2}{K_{gg}(\xi)}\,,\qquad K_{gg}(\xi)=\xi(1-\xi)-2+\frac{1}{\xi(1-\xi)}=\frac{\mathcal P_{gg+}(\xi)}{2N_c}\,,
\end{align}
(\texttt{b1}~[2], 15 digits). So a quark loop would add to the kernel $K_{gg}$ of the self-energy (Appendix~\dref{app:wavefunctions}) the kernel $K_q=\frac{n_fT_R}{N_c}\big[1-\frac{4\xi(1-\xi)}{d_\perp}\big]$ (the $d_\perp$-dimensional spin trace, $\frac{1}{d_\perp}\mathrm{tr}$ of the angular average), with $\int_0^1K_q=\frac{n_fT_R}{N_c}(\frac23-\frac\epsilon6)$ and $\int_0^1(\xi^2+(1-\xi)^2)\ln(\xi(1-\xi))=-\frac{13}{9}$ (\texttt{b2}~[1,2]). Also checked with the paper's own matrix elements (\texttt{b1}~[4]): the instantaneous emission of the pair, Eq.~(\dref{eq:m4}), cancels exactly the $\epsilon^-$ (longitudinal) part of the emission $\rho\to g\to gg$, Eqs.~(\dref{eq:m1}), (\dref{eq:m3}); the same holds for $q\bar q$ (the $\epsilon^-$ part of $\bar u\slashed\epsilon v$ is $Q^iS/Q^+$ with $S=\bar u\gamma^+v=2\sqrt{p_1^+p_2^+}\,\delta_{h,-\bar h}$, \texttt{b1}~[3]). This fixes the relative normalisation of all quark-loop pieces below without further input.

\paragraph{Where a quark loop would enter at this order.} (i) Row~I.II: the pole coefficient $\frac{11}{3}=-2\int K_{gg}$ (non-soft part) becomes $\frac{11}{3}-\frac43\frac{n_fT_R}{N_c}=\frac{\beta_0}{N_c}$, $\beta_0=\frac{11}{3}N_c-\frac43T_Rn_f$, and the constant becomes $\frac{67}{9}-\frac{\pi^2}{3}-\frac{20}{9}\frac{n_fT_R}{N_c}$ (the $-\frac{20}{9}$ is $-2(-\frac{13}{9}+\frac13)$, the $\frac13$ from the $O(\epsilon)$ part of $K_q$ times the pole), i.e.\ $\frac{2}{N_c}K_{\rm CMW}$, $K_{\rm CMW}=N_c(\frac{67}{18}-\frac{\pi^2}{6})-\frac{10}{9}T_Rn_f$ (\texttt{b2}~[2]). (ii) Row~I.V: a $q\bar q$ pair crossing the shock wave, colour $T_R\,U^{af}(\v y)-\mathrm{tr}\big(t^aU_F(\v x)t^fU_F^\dagger(\v z)\big)$, which reduces to $T_R[U(\v y)-U(\v x)]$ for a point-like pair, just as $N_cU(\v y)-f^{adc}f^{fbe}U^{db}(\v x)U^{ce}(\v z)\to N_c[U(\v y)-U(\v x)]$; its point-like UV part has $2L-\frac{11}{6}\to2L-\frac{\beta_0}{2N_c}$ ($q\bar q$ has no soft end points, no $L$). (iii) The normalisation of the dressed observed gluon: $\frac{11}{12}\to\frac{\beta_0}{4N_c}$ in both halves. (iv) The $\delta(1-\xi)$ of the $\MS$ $\mathcal P_{gg}$: $\frac{11N_c}{6}\to\frac{\beta_0}{2}$, and therefore the RC term. (v) There is no real $N_f$ term for a measured \emph{gluon} at this order (a $q\bar q$ final state contains no gluon); for a measured \emph{hadron} the quark-fragmentation channel would be an additional set of rows. The pole table with a general $\beta_0$ sums to $-\frac{\ell_\vee}{2}$ (single $\epsilon$) and to UV $\frac{\ell_\vee}{2}$, IR $-\ell_\vee$ (two $\epsilon$), independently of $N_f$ (\texttt{b2}~[4]).

\paragraph{The identification.}
\begin{center}\footnotesize
\begin{tabular}{@{}p{4.1cm}p{3.0cm}p{8.6cm}@{}}
\toprule
number in the paper & with a quark loop & what it is\\\midrule
Row I.II, $\frac{11}{3}$ (Eq.~\dref{eq:A3onerho}) & $\beta_0/N_c$ & UV coefficient of the gluon self-energy, $-2\int d\xi\,[K_{gg}+K_q]$ (non-soft part): the light-cone-gauge $Z_3$ of the emitted gluon\\
Row I.V, $2L-\frac{11}{6}$ ($N_{2b\,\rm UV}$, Eq.~\dref{eq:dg_V}) & $2L-\frac{\beta_0}{2N_c}$ & the same loop straddling the shock wave, pair unresolved; same $\xi$ integral\\
$\frac{11}{6}$ in $\mathfrak K^{\rm V}_{2b}$ (Eqs.~\dref{eq:sGIkernels}, \dref{eq:app_V2}) & $\frac{\beta_0}{2N_c}$ & integrand of the subtraction term of $N_{2b\,\rm UV}$ restored in $N_{2b\,\rm rest}$, $-\int d\xi[\xi(1-\xi)-2]$; a subtraction, not a running-coupling term (Sec.~\ref{sec:b-err})\\
normalisation, counterterm, RC: $\frac{11}{12}$ & $\frac{\beta_0}{4N_c}$ & $Z_3$ (UV half), coupling counterterm, and the virtual $\frac{\beta_0}{2}\delta(1-\xi)$ (IR half) completing $\mathcal P_{gg+}$ to the $\MS$ $\mathcal P_{gg}$\\
Rows I.II/I.V cancellation of the $\frac{11}{3}$ UV terms & $\beta_0$ on both sides & cancels for any $N_f$ (identity of Sec.~\ref{sec:b-group})\\
tail coefficient $-\frac{11}{6}$ (Eq.~\dref{eq:dg_tailtot}) & $-\frac{11}{6}$ & $\int\frac{d\xi}{\xi_+(1-\xi)_+}N(\xi)$, double-plus integral of the real $\mathcal P_{gg+}$ (\texttt{b2}~[1]); the $\delta(1-\xi)$ parts of DGLAP ($\frac{\beta_0}{2}$) and RC ($-\frac{\beta_0}{2}$) cancel in the tail for any $\beta_0$; not $\beta_0$\\
\bottomrule\end{tabular}\end{center}
In pure glue the coincidence $\frac{\beta_0}{2N_c}=-\int_0^1d\xi[\xi(1-\xi)-2]=\frac{11}{6}=-\int\frac{d\xi}{\xi_+(1-\xi)_+}N(\xi)$ is exact; only the quark loop separates the first equality (virtual, complete kernel) from the second (real $gg$ only). This is the content of the referee's first bullet of Major~3, and it is answered without computing the full $N_f$ rows.

\noindent\emph{Drop-in (end of Section~\dref{sec:dg_poles}):} ``\textbf{Remark (pure Yang--Mills).} The calculation is done for $N_f=0$, where $\beta_0/(2N_c)=\frac{11}{6}$ coincides with integrals of the splitting function. The identification of the $\beta_0$ terms can be made by the BLM device of adding a quark loop, without computing the $N_f$ terms of the result: a $q\bar q$ loop adds to the kernel $\xi(1-\xi)-2+\frac{1}{\xi(1-\xi)}$ of the self-energy the kernel $\frac{n_fT_R}{N_c}[\xi^2+(1-\xi)^2]$ (with the same normalisation of the vertices), and the corresponding $q\bar q$ pair crosses the shock wave in Row~I.V. Hence the $\frac{11}{3}$ of Row~I.II, the $-\frac{11}{6}$ of $N_{2b\,\rm UV}$, the $\frac{11}{6}$ of the subtraction term in $\mathfrak K^{\rm V}_{2b}$, the $\frac{11}{12}$ of the normalisation, of the coupling counterterm and of the RC term, and the $\frac{\beta_0}{2}$ of $\delta(1-\xi)$ become $\beta_0$, whereas the coefficient $-\frac{11}{6}$ of the large-distance tail, Eq.~\eqref{eq:dg_tailtot}, is the double-plus integral of the real $\mathcal P_{gg+}$ and stays $-\frac{11}{6}$. The constant of Row~I.II would become $\frac{67}{9}-\frac{\pi^2}{3}-\frac{20}{9}\frac{n_fT_R}{N_c}=\frac{2}{N_c}K_{\rm CMW}$.''

\subsection{The referee's error item on the $\frac{11}{6}$ term}\label{sec:b-err}
Row~I.V splits as $N_{2b}=N_{2b\,\rm UV}+N_{2b\,\rm rest}$ (Group\_I\_Row\_V, checked): $N_{2b\,\rm UV}$ is the integrand $\Phi_{\rm UV}=\frac{1}{k^+}K_{gg}(\xi)\,\Kv(x'-w')\cdot\Kv(y-x)\,\frac{(y-z)\cdot(y-x)}{(y-z)^2(x-z)^2}$ with colour $N_c[U(y)-U(x)]$; its $z$ integral in $d=2-\epsilon$ is $\pi[\frac2\epsilon+\mathcal B(y-x)]$, so its $-\frac{11}{6}$ part is Eq.~(\dref{eq:dg_V}), which cancels Row~I.II. $N_{2b\,\rm rest}=(\Phi-\Phi_{\rm UV})\times$colour contains $-\Phi_{\rm UV}$ at the integrand level; after the plus prescriptions its non-soft part is $-\frac{1}{k^+}[\xi(1-\xi)-2]\times(\dots)$, whose $\xi$ integral is the $+\frac{11}{6}$ of $\mathfrak K^{\rm V}_{2b}$. The two $\frac{11}{6}$ are the two halves of one add-and-subtract: the integrated subtraction cancels Row~I.II; the unintegrated one is needed in the finite term, where it makes $z\to x$ integrable (angular integral first, as the draft says after Eq.~\dref{eq:app_F}); its $z$ integral alone diverges. ``Cancels completely'' is true of the dimensionally regularised integral, not of the kernel. \emph{Drop-in (Section~\dref{sec:dg_poles}, after Eq.~\dref{eq:dg_V}):} ``The $-\frac{11}{6}$ of $N_{2b\,\rm UV}$ is the dimensionally regularised $z$ integral of a subtraction term; the same term, unintegrated and with the opposite sign, remains in the finite part of the row, where it is the term $\frac{11}{6}$ of $\mathfrak K^{\rm V}_{2b}$, Eq.~\eqref{eq:sGIkernels}, and makes the point $\v z\to\v x$ integrable. The cancellation against Row~I.II concerns the integrated subtraction only.''

\subsection{Rows I.II and I.V: what is a result and what is a grouping}\label{sec:b-group}
\paragraph{An exact identity.} Let $a=E_g(\v Q)$ be the energy of the emitted gluon (momentum $\v Q$), $D=\tilde{\v p}^2/(2\xi(1-\xi)Q^+)$ the relative energy of the pair and $b=a+D$ the energy of the pair. The self-energy part of $\mf A^{(3)}$ is $\psi_g\sum|V|^2/(ab)$; the part of Row~I.V in which the pair produced by the triple-gluon vertex (with the instantaneous term, which cancels the $\epsilon^-$ part, \texttt{b1}~[4]) is recombined by $\bar{\mf C}_1^\dagger$ without resolving the target is $\psi_g\sum|V|^2/(bD)$, and
\begin{align}\label{b:id}
\frac{1}{ab}+\frac{1}{bD}=\frac{1}{aD}\,,\qquad\int d^d\tilde p\,\frac{\sum|V|^2}{aD}=0\ \ (\text{scaleless})
\end{align}
(\texttt{b2}~[3]); the remaining $\epsilon^-$-type piece $S\,V^*$ is odd in $\tilde{\v p}$ and integrates to zero. The identity holds in $d$ dimensions for each $\xi$, with the same supports $\Lambda<p^+<k^+-\Lambda$ in both rows, and for the term with the rotated charge ($U(\v y)$) as well as for the term with the rotated pair. Hence, if the colour of the pair is evaluated at its \emph{centre of mass} $\v b=\xi\v x+(1-\xi)\v z$ (where the phase $e^{-i\v k\cdot(\v b-\v w')}$ already sits),
\begin{align}\label{b:split}
\text{Row I.II}\big|_{\rm self\text{-}energy}+\text{Row I.V}\big|_{\rm 3g,\ unresolved}=0\quad\text{exactly},\qquad
\text{Row I.V}\big|_{\rm 3g}=\int\Phi\,\big[N_cU^{af}(\v b)-f^{adc}f^{fbe}U^{db}(\v x)U^{ce}(\v z)\big]\,(\dots)\,,
\end{align}
pole, $\mathcal B$ logarithm, the soft end points of the self-energy and the constant $\frac{67}{9}-\frac{\pi^2}{3}$ included. (The commutator part of $\mf B_2$ in Row~I.V pairs with the commutator diagrams of $\mf A^{(3)}$ and is not covered by Eq.~\eqref{b:id}.)

\paragraph{Results.} (R1) The UV poles of Rows~I.II and I.V cancel, for any $N_f$ and in any polarisation scheme. (R2) $\ln\mu^2$ cancels between them. (R3) What remains of the $\beta_0$-type coefficient is in the resolving term of Eq.~\eqref{b:split}: since $\Phi\to K_{gg}(\xi)\times[\text{UV kernel}\propto1/(\v x-\v z)^2]$ for $|\v x-\v z|\ll|\v y-\v x|$ (checked in the project: ``the $z\to x$ behaviour of the full $2b$ integrand is reproduced by the UV term in any $d_\perp$''), the resolving term contains, for a target with saturation scale $Q_s$ and an emitter distance $|\v y-\v x|\gg1/Q_s$, the window $1/Q_s\ll|\v x-\v z|\ll|\v y-\v x|$ with coefficient $\int d\xi\,K_{gg}=2L-\frac{11}{6}$ ($\to2L-\frac{\beta_0}{2N_c}$) times $\pi\ln[(\v y-\v x)^2Q_s^2]$ times a colour factor that does not vanish there. So the coupling of the emission vertex does not run with the emitter distance only as long as the pair cannot resolve the target within that distance; beyond it a $\beta_0\ln[(\v x-\v w)^2Q_s^2]$-type term is present in the finite terms ($N_{2a}$, $N_{2b\,\rm rest}$) after the relabelling. This is a structural statement (power counting of the kernel); it was not evaluated with a target model.

\paragraph{Grouping.} The draft's subtraction evaluates the colour at $\v x$ (the gluon with $\xi$), with $\mathcal B(\v y-\v x)$. Any other point ($\v b$, $\v z$) or argument of $\mathcal B$ moves finite terms between ``Row I.II $+$ $N_{2b\,\rm UV}$'' and the finite kernels; the statement ``there is no finite term $\frac{\alpha_s\beta_0}{4\pi}[\mathcal B(\v x-\v w)+\dots]\times$LO'' is true for this grouping. \emph{Drop-in (replacing, in Section~\dref{sec:dg_poles}, ``Consequently the emission coupling does not run \dots\ survive~\cite{Kovner:2023vsy}.''):} ``The cancellation of the poles and of the dependence on $\mu$ between Rows~I.II and~I.V is a result. In the form used here, with the colour of the pair evaluated at the position $\v x$ of one of its gluons, no term $\frac{\alpha_s\beta_0}{4\pi}[\mathcal B(\v x-\v w)+\mathcal B(\v x'-\v w')]\times$LO remains explicitly; this depends on the grouping. The region in which the pair resolves the target is in the finite terms of Row~I.V, which carry the same coefficient $2L-\frac{11}{6}$ multiplying the small-pair kernel; for an emitter distance larger than the inverse saturation scale they contain a logarithm of $|\v x-\v w|Q_s$ with this coefficient. Hence the coupling of the emission vertex does not run with $\v x-\v w$ as long as the pair does not resolve the target within this distance.''

\paragraph{Consequence for the constant.}\label{sec:b-const} By Eq.~\eqref{b:split} the constant $\frac{\alpha_sN_c}{2\pi}(\frac{67}{9}-\frac{\pi^2}{3})\times$LO of Row~I.II, Eq.~(\dref{eq:constantLO}), is not a grouping-invariant part of the result: in the centre-of-mass grouping it is cancelled identically by Row~I.V; in the draft's grouping (colour at $\v x$) the compensating part sits in $N_{2b\,\rm rest}$, up to the terms generated by $U(\v b)-U(\v x)$. This is consistent with the condition found in the dilute check (\texttt{ref3\_dilute}, Eq.~\texttt{d:cond}: Group~I must contain no net $\frac{67}{9}-\frac{\pi^2}{3}$ beyond scheme terms) and contradicts the draft's reading of Eq.~(\dref{eq:constantLO}) as the LO-type constant of the result (Section~\dref{sec:lit_constant}). \textbf{Not done:} the value of the LO-type constant left in $N_{2b\,\rm rest}$ in the draft's grouping. Recommended: the sentence in Section~\dref{sec:lit_constant} ``the bare-vertex value~\eqref{eq:constantLO} is the appropriate one'' should be withdrawn until this is evaluated.

\subsection{Independent renormalisation and factorisation scales}\label{sec:b-scales}
The single-$\epsilon$ bookkeeping cancels the coupling counterterm against the collinear pole of the $\delta(1-\xi)$ term, which requires one $\mu$. With $g^2_0=\mu_R^\epsilon g^2(\mu_R)[1-\frac{\alpha_s\beta_0}{4\pi}\frac2\epsilon]$ and $D_0=[\delta(1-\xi)+\frac{\alpha_s}{2\pi}\frac2\epsilon\mathcal P_{gg}]\otimes D(\mu_F)$, the two halves of the scaleless normalisation are cancelled separately (UV half by the coupling counterterm at $\mu_R$, IR half by the $\delta(1-\xi)$ part of the FF counterterm at $\mu_F$), and the finite terms are
\begin{align}
\frac{d^3N_{\rm LO}}{d^3k}+\frac{d^3N_{\rm RC}}{d^3k}&=\frac{1}{(2\pi)^3}\frac{g^2(\mu_R)}{2\pi^2}\frac{1}{k^+}\int_{x,y,x',y'}e^{-ik\cdot(y-y')}\Big[1+\frac{\alpha_s\beta_0}{4\pi}\ln\frac{(y'-y)^2\mu_R^2}{c_0^2}\Big]\mathcal O(y,y')\,,\notag\\
\frac{d^3N_{\rm DGLAP}}{d^3k}&=-\mathcal C\int_0^1d\xi\,\widetilde{\mathcal P}_{gg}(\xi)\,\mathcal F(\xi;\mu_F)\,,\qquad \mathcal F(\xi;\mu_F)=\mathcal F(\xi)\big|_{\mu\to\mu_F}\,,
\end{align}
including the $\frac{\beta_0}{2}\delta(1-\xi)$ term at $\mu_F$. Their $\beta_0$ part is $N_{\rm RC}(\mu_R)-\mathcal C\frac{\beta_0}{2}\mathcal F(1;\mu_F)=+\frac{\alpha_s\beta_0}{4\pi}\ln\frac{\mu_R^2}{\mu_F^2}\times$LO, and the result is $\mu_R$- and $\mu_F$-independent to this order (\texttt{b2}~[5]: remainders $O(\alpha_s^3)$). In the local form, $N^{\rm loc}_{\rm RC}=-\frac{\alpha_s\beta_0}{4\pi}\ln\frac{\v k^2}{\mu_R^2}\mathcal N_0\widehat L(k)$ and $N^{\rm loc}_{\rm DGLAP}=\mathcal C\ln\frac{\v k^2}{\mu_F^2}\int_0^1d\xi\,\mathcal P_{gg}\Pi$; $\mathcal R$ and $N_{\mathcal R}$ are unchanged. The draft's statement that the residual $\mu$ dependence comes only from the $\delta(1-\xi)$ term holds only at $\mu_R=\mu_F$. \emph{Drop-in (after Eq.~\dref{eq:LORC}):} ``With separate scales the coupling is $g^2(\mu_R)$ in Eq.~\eqref{eq:LORC}, the DGLAP term~\eqref{eq:DGLAP} is evaluated at $\mu_F$, including its $\delta(1-\xi)$ term, and their sum contains $\frac{\alpha_s\beta_0}{4\pi}\ln\frac{\mu_R^2}{\mu_F^2}\times$LO; the add-and-subtract of the $\delta(1-\xi)$ term identifies $\mu_R$ with $\mu_F$.''

\subsection{Scheme}\label{sec:b-scheme}
\paragraph{What the paper does.} Polarisation sums $\delta_{ii}=d_\perp=2-\epsilon$ with $d_\perp$-dimensional momenta (CDR) wherever a pole multiplies them: Row~I.II, the reduction $\tilde p^l\tilde p^i\to\delta^{li}\tilde p^2/d$ with $d_\perp$ internal polarisations gives the factor $d_\perp/d=1$ in front of $\xi(1-\xi)$ (main file, self-energy; \texttt{b2}~[2]); Row~I.V, $N_{2b\,\rm UV}$: $d_\perp\cdot\frac{1}{d_\perp}=1$ (record, Divergent terms l.112); Row~II.IV: $d_\perp=2-\epsilon$ kept in the contraction (record, Group II Row IV l.86, l.143); the finite rows are locally integrable, so $d_\perp=2$ there (record, Group II l.81; Row~II.I $N_{2f}$). The real two-gluon emission has the CDR $(D-2)$ structure of Fadin--Lipatov/BSV for any $d_\perp$ (\texttt{ref3\_dilute}). The counterterms are pure poles: $\MS$ coupling (CDR) and $\MS$ fragmentation function with the four-dimensional $\mathcal P_{gg}$ (in CDR the averaged $\hat P_{gg}$ has no $O(\epsilon)$ term), with the $\xi^{-d}$ Jacobian of Section~\dref{sec:dg_msbar}. The constant $\frac{67}{9}-\frac{\pi^2}{3}$ and Eq.~(\dref{eq:A3onerho}) are CDR values. All pieces are in one scheme.

\paragraph{Two-dimensional polarisations, $d$-dimensional momenta (DRED-like).} The factor becomes $2/d=1+\frac\epsilon2+\dots$ (\texttt{b2}~[2]): Row~I.II, $\frac{67}{9}-\frac{\pi^2}{3}\to\frac{64}{9}-\frac{\pi^2}{3}$ ($K_{\rm CMW}\to K_{\rm CMW}-\frac{N_c}{6}$, i.e.\ $-\frac{\alpha_sN_c}{6\pi}\times$LO); Row~I.V, $N_{2b\,\rm UV}$ gets $+\frac16$ per side, $+\frac{\alpha_sN_c}{6\pi}\times$LO, which cancels the former (as it must by Eq.~\eqref{b:id}); Row~II.IV gets $-\mathcal C\int_0^1d\xi\,2N_c\,\xi(1-\xi)\,\Pi(\xi)$. With $\alpha_s^{\rm DR}=\alpha_s^{\MS}(1+\frac{\alpha_sN_c}{12\pi})$ and $D^{\rm DR}=D^{\MS}\otimes[\delta+\frac{\alpha_s}{2\pi}c]$, $c(\xi)=2N_c\xi(1-\xi)-\frac{N_c}{6}\delta(1-\xi)$, the result is unchanged; the $\delta$ coefficient $\frac{N_c}{6}$ is fixed independently by the momentum sum rule, $\int\xi c=0$, and equals the coupling conversion, $\frac{\alpha_s}{2\pi}\frac{N_c}{6}=\frac{\alpha_sN_c}{12\pi}$ (\texttt{b2}~[6]). The shift of order $N_c/6$ the referee mentions is this one; in the paper's combination of rows it cancels in the virtual part and appears only in the conversion of $\alpha_s$ and $D$.

\noindent\emph{Drop-in (Section~\dref{sec:dg_poles}, first paragraph):} ``We use conventional dimensional regularisation: transverse momenta and gluon polarisations are continued to $d_\perp=2-\epsilon$ dimensions wherever a polarisation sum multiplies a pole (Rows~I.II, I.V and II.IV); in the finite terms $d_\perp=2$. The coupling is the $\MS$ coupling and the fragmentation function the $\MS$ number density (pure-pole counterterms with the four-dimensional $\mathcal P_{gg}$), both in this scheme; Eq.~\eqref{eq:A3onerho} and the constant $\frac{67}{9}-\frac{\pi^2}{3}$ refer to it. With two-dimensional polarisations the constant of Row~I.II would be $\frac{64}{9}-\frac{\pi^2}{3}$, that of Row~I.V would change by the opposite amount, Row~II.IV by $-\mathcal C\int d\xi\,2N_c\xi(1-\xi)\Pi(\xi)$, and the coupling and fragmentation function would have to be converted, $\alpha_s^{\rm DR}=\alpha_s^{\MS}(1+\frac{\alpha_sN_c}{12\pi})$.''

\subsection{Not done}
Full coordinate-space quark-loop rows (not needed, $N_f=0$); the LO-type constant left in $N_{2b\,\rm rest}$ in the draft's grouping (Sec.~\ref{sec:b-const}); a target-model evaluation of the resolving window (R3); the sign of the Row~II.IV DRED term was obtained from the general $d_\perp/d$ rule (consistent with the record's ``$+1$'' in the bracket for $d_\perp=2$) and the closure of \texttt{b2}~[6], not by re-deriving the row.
